\section{Related Work}
\label{sec:related}

\sstitle{LLM and reference-based phishing detection}
The dominant modern detectors decide whether a page impersonates a brand it does
not own. Reference-based systems match a rendered page against a brand library of
logos and layouts~\cite{abdelnabi2020visualphishnet,lin2021phishpedia,
liu2022phishintention}, and recent work refreshes that reference set
online~\cite{dynaphish2023} or replaces it with an LLM that recognizes the brand
and checks the domain directly~\cite{tan2024phishllm,li2024knowphish}. A parallel
line treats the page as URL or image features for a learned
classifier~\cite{diffusion2025phishing,aljofey2026hybrid}, surveyed
broadly~\cite{slr2023phishing}. Across these sub-mechanisms the shared target is
\emph{aggregate} accuracy, reported as a single page-level verdict. Building on this
line rather than competing with it, \toolname{} invokes such detectors as verification
tools (\autoref{sec:ground}) -- inheriting their accuracy so that it stays
\emph{competitive} at full coverage (\autoref{tab:detect}) -- and adds what a bare
detector lacks: a \emph{per-verdict} decision about whether the verdict can be trusted,
acted on automatically, and explained (\autoref{sec:select}).

\sstitle{Multi-agent and ensemble phishing detection}
Closest to us are detectors that run several models. Majority voting over LLMs has
been studied directly for phishing~\cite{ensembleornot2024}; multi-agent systems
assign agents to different signals and either debate to a verdict or merge their
rationales into one explanation~\cite{phishdebate2025,multiphishguard2025},
sometimes backed by an episodic memory of past cases~\cite{memophish2026}.
Unlike these systems, which aggregate or consolidate at the level of the
\emph{label} or a fused rationale, \toolname{} measures agreement at the level of
individually cited evidence units (\eqref{eq:agree}) and treats \emph{disagreement
on evidence} as the signal to abstain rather than a conflict to resolve. The contrast is
sharpest with role-based debate frameworks such as PhishDebate~\cite{phishdebate2025},
which assign agents to fixed signal roles and \emph{resolve} their disagreement into a
single verdict through a moderator; \toolname{} instead leaves the unit of agreement at
the cue, never forces a resolution, and reads the residual disagreement as the very
quantity that gates trust -- so a page the debate would push to a confident verdict is one
\toolname{} declines to act on. The closest LLM-consolidator system,
MultiPhishGuard~\cite{multiphishguard2025}, has a separate model read the agents'
rationales and emit a deliberated verdict-and-confidence -- a stronger baseline than
label voting (B6 in \autoref{sec:experiments}), but one that still scores trust at the
label channel after a free-text fusion: a consolidator that finds the agents
``agree'' returns high confidence whether or not the agreed reasons hold on the page,
and the confidence is read directly without held-out calibration. \toolname{} differs
on three axes that the empirical comparison isolates: (i) the unit of agreement is a
typed, tool-checkable cue rather than a fused rationale; (ii) the score is
explicitly calibrated against held-out labels, so its reported probability of
correctness is directly actionable (the gap in ECE in \autoref{tab:main}); and (iii)
the act/abstain rule is set at a target selective risk rather than left to the
consolidator's self-reported confidence, which is what produces the fail-safe
behaviour under evidence-targeted attack (\autoref{sec:exp_adv}). Its agents
are tool-using investigators that each \emph{plan} their own investigation across
signals and modalities, and whose tool-grounded findings are \emph{aggregated}
(\autoref{sec:ground}) into a shared set of cited evidence cues that serves as both the
trust signal and the explanation, rather than a fused free-text rationale.

\sstitle{Label-level reliability and label-free evaluation}
A verdict's reliability is usually read from the label distribution: post-hoc
confidence calibration~\cite{guo2017calibration}, agreement across samples of one
model~\cite{wang2023selfconsistency}, or agreement across many models, as in
crowd-sourcing aggregation~\cite{dawid1979maximum}, agreement-on-the-line under
shift~\cite{baek2022agreement}, and trust calibration for phishing
detectors~\cite{tci2025}. Where these estimate reliability from \emph{how often
labels concur}, \toolname{} estimates it from \emph{whether the reasons concur and
hold}, which separates pages on which the panel agrees by coincidence from those on
which it agrees because the page truly carries the evidence (\eqref{eq:gea}).

\sstitle{Explanation faithfulness}
Causal audits show that an LLM's stated rationale is often not the driver of its
decision~\cite{ariadne2026}, and multi-agent debate has been used to \emph{improve}
the faithfulness of generated explanations~\cite{madr2024}. Such audits operate on
a single model and require counterfactual re-execution. In contrast, \toolname{}
uses cross-model agreement on grounded evidence as an inexpensive faithfulness
\emph{proxy}: a reason that several independent models cite and that the page
confirms (\eqref{eq:ground}) is, by construction, less likely to be a post-hoc
rationalization, at no extra re-execution cost.

\sstitle{Selective prediction}
We build on selective classification, which trades coverage for accuracy via an
abstention rule~\cite{geifman2017selective}, and on distribution-free uncertainty
quantification~\cite{vovk2005algorithmic}. Building on this foundation,
\toolname{} supplies the missing ingredient for the phishing setting: an
abstention score derived from grounded cross-model agreement (\eqref{eq:gea}),
calibrated on a held-out labeled split so the risk--coverage operating point is set
against real outcomes (\autoref{sec:select}).

\sstitle{Position and scope}
\toolname{} contributes a single new object, the Grounded Evidence-Agreement
score, and the selective detector built on it. \autoref{tab:related} positions it
against one representative per family along five capabilities -- evidence-level
agreement, page-grounding, calibrated trust, selective abstention, and a
verifiable explanation as output -- and \toolname{} is the only entry that
supports all five.
\begin{table}[!h]
\caption{Capability comparison for per-verdict trust in multi-model phishing
detection. \CIRC: yes; \HCIR: partial; \OCIR: no. Columns are defined in the body;
full-coverage detection accuracy is reported separately in \autoref{tab:detect}.}
\label{tab:related}
\footnotesize
\setlength{\tabcolsep}{2pt}
\begin{tabular*}{\columnwidth}{@{\extracolsep{\fill}}lccccc@{}}
\toprule
 & Evid.- & Page- & Calib.- & Selective & Explan. \\
Approach & level & grounded & trust & abstain & output \\
\midrule
Phishpedia~\cite{lin2021phishpedia}            & \OCIR & \HCIR & \OCIR & \OCIR & \HCIR \\
PhishLLM~\cite{tan2024phishllm}                & \OCIR & \HCIR & \OCIR & \OCIR & \HCIR \\
Single conf.~\cite{guo2017calibration}        & \OCIR & \OCIR & \HCIR & \CIRC & \OCIR \\
Self-consist.~\cite{wang2023selfconsistency}   & \OCIR & \OCIR & \OCIR & \HCIR & \OCIR \\
Majority vote~\cite{ensembleornot2024}         & \OCIR & \OCIR & \OCIR & \HCIR & \OCIR \\
Label agree.~\cite{dawid1979maximum}           & \OCIR & \OCIR & \HCIR & \HCIR & \OCIR \\
MultiPhishGuard~\cite{multiphishguard2025}     & \HCIR & \OCIR & \OCIR & \HCIR & \HCIR \\
Faithfulness audit~\cite{ariadne2026}          & \HCIR & \OCIR & \OCIR & \OCIR & \HCIR \\
\textbf{\toolname{}} (ours)                          & \CIRC & \CIRC & \CIRC & \CIRC & \CIRC \\
\bottomrule
\end{tabular*}
\end{table}

We focus on per-verdict trust and explanation for a multi-model
phishing detector; the framework is agnostic to the specific base models and
evidence schema. We do \emph{not} claim (i) a new base
detector or a leaderboard win on full-coverage accuracy -- \toolname{} reuses existing
detectors as tools and is \emph{competitive} with them (\autoref{tab:detect}),
contributing the per-verdict trust and explanation layer on top; (ii) robustness to an
adversary who controls the panel agents or corrupts the tools' inputs, a
harness-level threat we treat as out of scope; or (iii) coverage guarantees on
fully cloaked pages that carry no verifiable cues, where the
detector necessarily abstains rather than errs.
